\documentclass{article}
\usepackage{geometry} % Required for customizing margins
\geometry{margin=1in} % Sets 1-inch margins on all sides
\usepackage{listings}
\usepackage{tikz}
\usepackage{titlesec}  % For customizing section/subsection titles

% Custom format for sections and subsections
\titleformat{\section}[block]{\normalfont\Large\bfseries}{\thesection}{1em}{}
\titleformat{\subsection}[block]{\normalfont\large\bfseries}{\thesubsection}{1em}{}

% Page layout settings
\geometry{
    top=1.0in,
    bottom=1.0in,
    left=1.0in,
    right=1.0in
}

\begin{document}

% Horizontal line
\hrule
\vspace{0.2cm}

% Centered title
\begin{center}
    \LARGE \textbf{CSCE 350} \\
    \Large Data Structures and Algorithms \\
    \large HW \#1
\end{center}

% Horizontal line
\hrule
\vspace{0.2cm}

% Author and semester info
\begin{center}
    \textbf{Author:} Jacob Vaught \\
    \vspace{0.2cm}
    \textbf{Semester:} Fall 2023 \\
\end{center}

% Another horizontal line
\hrule
\vspace{0.5cm}

\newpage
\section*{Problem \#1}
\newcommand{\drawstack}[2]{%
  \begin{minipage}[b]{0.1\textwidth}
    \centering
    \begin{tikzpicture}[baseline={(current bounding box.south)}]
      \ifx&#1&
        \draw (0,0) rectangle ++(1,1);
        \node at (0.5,0.5) {null};
      \else
        \foreach \x [count=\xi] in {#1}
        {
          \draw (0,\xi-1) rectangle ++(1,1);
          \node at (0.5,\xi-0.5) {\x};
        }
      \fi
    \end{tikzpicture} \\
    #2
  \end{minipage}%
  \vspace{0.05cm}% <-- Adjust this value to change vertical spacing
}

Assume you have an empty stack, a sequence of operations are
performed on it as shown below. Show all intermediate steps and the final stack.
\newline Operations: push(a), pop(), push(b), push(c), pop, push(d), pop, push(e), pop, pop
\newline \newline
\noindent
\drawstack{}{Initial Stack}
\drawstack{a}{After push(a)}
\drawstack{}{After pop()}
\drawstack{b}{After push(b)}
\drawstack{b,c}{After push(c)}
\drawstack{b}{After pop()}
\drawstack{b,d}{After push(d)}
\drawstack{b}{After pop()}
\drawstack{b,e}{After push(e)}
\drawstack{b}{After pop()}
\drawstack{}{Final Stack}

\newpage
\section*{Problem \#2}
\newcommand{\drawqueue}[2]{%
  \begin{minipage}[b]{0.25\textwidth}
    \centering
    \begin{tikzpicture}[baseline={(current bounding box.south)}]
      \ifx&#1&
        \draw (0,0) circle [radius=0.5];
        \node at (0,0) {0};
      \else
        \foreach \x [count=\xi] in {#1}
        {
          \draw (\xi,0) circle [radius=0.5];
          \node at (\xi,0) {\x};
        }
      \fi
    \end{tikzpicture} \\
    #2
  \end{minipage}%
  \vspace{0.05cm}% <-- Adjust this value to change vertical spacing
}



Assume you have an empty queue, and a sequence of operations are performed on it.
\newline\newline
\noindent
\drawqueue{}{Initial Queue}
\drawqueue{a}{After enqueue(a)}
\drawqueue{a,b}{After enqueue(b)}
\drawqueue{b}{After dequeue}
\drawqueue{b,c}{After enqueue(c)}
\drawqueue{b,c,d}{After enqueue(d)}
\drawqueue{b,c,d,e}{After enqueue(e)}
\drawqueue{c,d,e}{After dequeue}
\drawqueue{d,e}{After dequeue}
\drawqueue{d,e,f}{After enqueue(f)}

\newpage
\section*{Question 3}

\subsection*{Solution Description}
For solving the locker toggling problem, I followed my preferred development method: I initially implemented the algorithm in Python, which is my go-to language for prototyping. Once I verified that the algorithm worked as intended, I proceeded to convert it into C++. The following C++ code snippet illustrates the core logic:

\begin{lstlisting}[language=C++, basicstyle=\ttfamily]
std::vector<std::string> lockers(n + 1, "closed");
for(int i = 1; i <= r; ++i) {
    for(int j = i; j <= n; j += i) {
        if(lockers[j] == "closed") {
            lockers[j] = "open";
        } else {
            lockers[j] = "closed";
        }
    }
}
\end{lstlisting}
This C++ implementation uses a vector of strings to track each locker's status, as either 'open' or 'closed'. The two nested for-loops are responsible for the toggling operation according to the defined rules.

\subsection*{Challenges and Solutions}
The initial challenge was in deciding the most efficient way to toggle the locker status. Initially, I thought about using integers (0 for 'closed' and 1 for 'open'), but I opted for strings for better readability and code maintenance. Since I first implemented the solution in Python, transitioning to C++ was rather smooth, and any challenges in the logic were easily tackled using conditional statements, as you can see in the above code snippet.

\subsection*{Special Compilation Instructions}
To compile the C++ code, first save it as \texttt{LockerDoors\_Vaught\_Jacob.cpp}. Then, use the following shell commands:

\begin{verbatim}
g++ LockerDoors_Vaught_Jacob.cpp -o LockerDoors_Vaught_Jacob
./LockerDoors_Vaught_Jacob
\end{verbatim}
Executing these commands will create an executable named \texttt{LockerDoors\_Vaught\_Jacob}. Running this executable will prompt the user to input the number of lockers and the number of passes. The program will then display the final status of each locker.

\end{document}
